% LaTeX table fragment. Requires \usepackage[utf8]{inputenc}, \usepackage{booktabs,longtable,array}, and \usepackage{pdflscape}.
\begingroup
\scriptsize
\setlength{\tabcolsep}{3pt}
\begin{landscape}
\begin{longtable}{>{\raggedright\arraybackslash}p{4.8cm} r r r}
\caption{Phase B: H4: Rohstoffanteil der IMF-Konditionalität — Modellvergleich (1992–2023)}\label{tab:phase_b_h4_modellvergleich}\\
\toprule
Kennzahl & Basis-FE & Vollmodell-FE & Länder- und Jahres-FE
 \\
\midrule
\endfirsthead
\toprule
Kennzahl & Basis-FE & Vollmodell-FE & Länder- und Jahres-FE
 \\
\midrule
\endhead
\midrule
\multicolumn{4}{r}{Fortsetzung auf der naechsten Seite}\\
\endfoot
\bottomrule
\endlastfoot
Fokuseffekt & unsc3:resource\_dep & unsc3:resource\_dep & unsc3:resource\_dep \\
Koeffizient & 6.61e-04 & 5.01e-05 & -1.93e-06 \\
Robuster Standardfehler & 5.00e-04 & 7.02e-04 & 4.96e-04 \\
p-Wert & 0.187 & 0.943 & 0.997 \\
Beobachtungen & 380 & 212 & 212 \\
Länder & 104 & 62 & 62 \\
Länder mit unsc3 = 1 & 27 & 16 & 16 \\
Jahre & 1992--2023 & 1992--2023 & 1992--2023 \\
Länder-FE & Ja & Ja & Ja \\
Jahres-FE & Nein & Nein & Ja \\
Kontrollen & Laufzeit & DSV-Vollsatz & DSV-Vollsatz \\
\end{longtable}
\noindent\footnotesize Koeffizienten der Hypothese, nicht die Koeffizienten aller Kontrollen. Heteroskedastizitaetskonsistente robuste Standardfehler. Signifikanz: * p<0.10, ** p<0.05, *** p<0.01.\par
\end{landscape}
\endgroup

